%! Author = charon
%! Date = 6/4/24

\subsection{Desocketing von Software}\label{subsec:desocketing-von-software}
Das Desocketing von Software ist ein Prozess, bei dem die Kommunikation zwischen zwei oder mehreren Software-Modulen,
die über Sockets miteinander verbunden sind, entfernt wird.
Das Ziel des Desocketing-Prozesses ist es, die Software so zu modifizieren, dass sie ohne die Verwendung von Sockets funktioniert.
Dies kann aus verschiedenen Gründen erforderlich sein, z.B.\ wenn die Software auf einem System ohne Netzwerkverbindung ausgeführt werden soll
oder wenn die Software auf einem System ausgeführt werden soll, auf dem die Verwendung von Sockets aus Sicherheitsgründen eingeschränkt ist.
In der Regel wird der Desocketing-Prozess durchgeführt, indem die Software so modifiziert wird, dass sie die Sockets durch eine andere
Kommunikationsmethode ersetzt.
Dies kann z.B.\ durch die Verwendung von Pipes, Shared Memory oder anderen Inter-Process Communication (IPC) Mechanismen erfolgen. \newline
Im Fall dieser Arbeit wird eine Bibliothek implementiert, welche die Kommunikation des \gls{zup} über Sockets durch eine
die Kommunikation eines Streams ersetzt.
Diese Bibliothek wird in die Firmware des \gls{zup} integriert und ermöglicht es, die Kommunikation des \gls{zup} über
einen Stream zu steuern.
Dadurch wird es möglich mit der von \gls{afl} bereitgestellten \gls{tcp}-Fuzzing-Infrastruktur die Netzwerkkommunikation
des \gls{zup} auf Schwachstellen zu testen.
Dieser Schritt ist unter Verwendung von \gls{afl} erforderlich, da \gls{afl} nur Eingaben an ein Programm -- in Form
von Steams mittels Dateien -- übergeben kann.\newline

\noindent Zur Implementierung des Desocketing-Prozesses wird die Socket-API des \gls{zup} analysiert und die Kommunikation über Sockets
durch eine Kommunikation über Streams ersetzt.
